\chapter{Conclusion}
\label{ch:conclusion}

This thesis asked how an endpoint on a path-aware Internet should choose the best use of its paths given
the intention of the application and the state of the network. It answered in two movements that share one
idea, drawn from Mortise~\cite{Shen2026}: treat application intent as a first-class input that reshapes a
network mechanism, not as a constant compiled into it.

\Cref{ch:single_path} turned a fixed-objective DQN path selector for SCION~\cite{Bourak2025} into an
intent-conditioned one. Two design decisions carried the work: a permutation-equivariant \emph{scoring}
architecture that handles a variable, changing set of candidate paths, and the routing of the intent vector
into the \emph{per-path comparison} itself, which lets a single policy serve a range of objectives --
including ones it never trained on -- without retraining. \Cref{ch:multipath} asked what happens if we use several paths at once and optimize the
application itself, and built a
two-timescale hierarchical agent pair that co-optimizes video bitrate and per-path traffic split over
Multipath QUIC, coupled so that each agent accounts for the other.

\paragraph{What we learned.} The most useful findings were the ones that cut against intuition. First, that
\emph{where} intent is injected matters more than \emph{that} it is injected: for a comparison-based
selector, a conditioning signal reaching only the terms shared by every candidate leaves decisions
unchanged, because it cancels under the argmax; the intent has to reach each alternative's own features to
reorder them. Second, that the value of learned multipath scheduling is
\emph{not} a fixed quantity but a function of network volatility -- negligible when paths are stable, where a
good bitrate controller over a trivial split suffices, yet decisive under churn and bursts, where the same
app-only configuration collapses. Both findings sharpen when, and when not, to reach for the more
sophisticated mechanism.

\paragraph{Future work.} Several threads remain open, and much of the draft's structure is deliberately laid
out to guide them. On the experimental side: adding seed-level statistics and the flat-versus-scoring
comparison for both systems, testing the single-path selector on intents and source--destination pairs it
never trained on (\cref{sec:p1eval:zeroshot,sec:p1eval:dynamic}), and porting the non-stationary dynamics
into the NS-3 backend so packet-level results can confirm the mock findings. On the system side: replacing the TCP-like subflows with
a real Multipath QUIC transport, and enriching the VMAF surrogate so quality responds to delivery
conditions, not just bitrate. On the research side, three directions stand out: \emph{inferring} intent from
traffic rather than assuming it is declared; \emph{online and safe} learning against live traffic with
bounded worst-case behavior; and pushing the ``one principle, every controllable decision'' recipe further
-- for instance, conditioning the multipath agents on an explicit intent vector as the single-path selector
does, unifying the two systems into a single intent-conditioned, cross-layer controller for path-aware
networks.
